\chapter{Discussão e Conclusão}\label{ch:discussion}

\section{Interpretação dos resultados}\label{sec:interpretation}

O zoneamento responde à QP1: o agrupamento descorrelacionado divide GO~e~DF em quatro zonas
biofisicamente distintas. São elas: um planalto de cultivo alto e argiloso; uma escarpa íngreme, onde a
aptidão das culturas colapsa e a conservação é também o melhor uso absoluto; uma baixada drenada,
voltada à piscicultura; e uma baixada quente e arenosa, onde o calor deprime mais as culturas do que a
pastagem (Seção~\ref{sec:zoning-results}, Tabela~\ref{tab:zone-suit}).

A leitura que essas quatro zonas sustentam é, porém, mais contida do que quatro rótulos sugerem.
Só duas delas têm vocação comparativa realmente definida (a escarpa de conservação e a baixada
drenada). Nas outras duas, que somam três quartos do território, a margem sobre a vocação
secundária é uma ordem de grandeza menor, e a zona deve ser lida pelo par de usos. Com $k = 4$ e sete
segmentos, três segmentos não recebem rótulo primário em nenhuma zona. A soja, embora tenha aptidão
absoluta alta em quase todas as zonas de cultivo, nunca é o melhor uso comparativo. A plausibilidade
da partição, portanto, não vem da nitidez dos rótulos. Vem do fato de que, embora construída apenas
com atributos biofísicos, ela separa fortemente a produtividade observada.

O padrão subjacente é que o território de GO~e~DF é climaticamente quase uniforme: quem separa terra
boa de terra ruim é o solo e o terreno, e não o clima. Por isso, essa leitura depende de dois produtos
globais (solo e terreno) não validados localmente. É a acurácia dessas camadas, e não a do clima, que
se propaga para as fronteiras entre zonas (Seção~\ref{sec:limitations}).

A projeção climática estende essa leitura a meados do século (QP3). A mudança é de magnitude
pequena a moderada, mas com sinal coerente em todo o conjunto de modelos. Para as culturas, ela
acentua o gradiente já existente em vez de redesenhá-lo: as perdas recaem sobre a baixada quente
e arenosa, já marginal, enquanto o planalto de cultivo permanece essencialmente inalterado. Três
segmentos são os mais expostos, e a ordem entre eles inverte o que uma leitura ingênua sugeriria. O
mais afetado não é uma cultura: é a \textbf{geração fotovoltaica}, que, sob SSP5-8.5 em 2051--2070,
perde mais da metade da área apta pela redução de desempenho térmico dos painéis, um efeito direto do 
aumento da temperatura. Em seguida vem a conservação. Isso é coerente com o fato de a
vegetação nativa não ter manejo que a proteja do clima, ao contrário das culturas, que o manejo mantém
perto do ótimo térmico. Só então aparecem as outras culturas anuais, com a maior perda de área apta
entre as culturas alimentares. Soja, cana-de-açúcar e pecuária permanecem praticamente intactas, e a
piscicultura é o único segmento que ganha área apta, porque o aquecimento empurra a temperatura do
trimestre mais quente \emph{em direção} ao seu ótimo, e não para fora dele (Seção~\ref{sec:shift},
Tabela~\ref{tab:shift-share}).

A análise do uso atual responde à QP2 e corrobora a QP1: na metade do território reservada à
validação, a discriminação por presença é forte para a cana-de-açúcar e a soja
(Seção~\ref{sec:validation-present}). Duas tensões merecem registro. Primeira: uma lacuna substancial,
majoritariamente sobre pastagem, separa potencial e uso atual. Cerca de dois quintos do território são
aptos à soja e ainda não cultivados com ela, e a lacuna das outras culturas anuais é a maior entre os
quatro segmentos para os quais ela é calculada (Tabela~\ref{tab:pot-real-gap}). Isso aponta para
intensificação, e não para abertura de novas áreas, o que é coerente com a trajetória já documentada
para a fronteira agrícola do Cerrado
\citep{spera_agricultural_2017,caballero_transformation_2023}. Segunda: a verificação cruzada
orientada por dados mostra que o mapa baseado em conhecimento é mais permissivo que o nicho realizado,
com concordância apenas razoável. Um $\kappa$ inferior ao que a AUC sugeriria não é contradição: as
duas métricas respondem a perguntas diferentes sobre o mesmo resultado
(Seção~\ref{sec:validation-present}).

A leitura econômica implícita na primeira dessas tensões tem respaldo na literatura sobre a
fronteira agrícola do Cerrado. Em Goiás, entre 2003 e 2016, o aumento da produção veio sobretudo da
intensificação, e não da abertura de novas áreas. A lotação do rebanho bovino subiu de 1,17 para
1,24 cabeças por hectare, a área de dupla safra soja-milho passou de cerca de 373~mil para 1,4~milhão
de hectares e, ainda assim, entre 50\% e 60\% das pastagens do Cerrado seguem degradadas
\citep[pp.~2, 4]{spera_agricultural_2017}. Há, portanto, espaço físico para intensificar sem converter
vegetação nativa, que é exatamente a leitura sugerida pela lacuna entre potencial e uso atual deste
trabalho.

Esse espaço, porém, não é gratuito. Recuperar pastagem degradada para lavoura ou pecuária mais
intensiva tem um custo de entrada que a aptidão biofísica, por si só, não captura: calagem e
correção da acidez do solo, mecanização e acesso a crédito. O crédito, por sua vez, é limitado pela
insegurança da posse da terra, que afeta parte expressiva dos pequenos produtores brasileiros. Em Goiás,
porém, mais de 80\% da área agrícola estava, em 2006, em propriedades acima de 100~ha
\citep[pp.~2, 5]{spera_agricultural_2017}. Em escala de bioma, a conversão de pastagem para soja e
outras culturas está associada a fatores econômicos (preços de exportação e demanda por alimentos e
biocombustíveis), e não apenas à disponibilidade de terra apta
\citep[pp.~10, 12]{caballero_transformation_2023}. Preço, custo de recuperação do solo, posse da terra
e logística de escoamento são, portanto, candidatos plausíveis para explicar a lacuna de 41,5\%
observada para a soja, mas nenhum deles é modelado neste trabalho. Testá-los é o objeto da camada
econômica proposta na Seção~\ref{sec:futurework}.

Algumas escolhas de delineamento merecem registro à parte dos mapas, e cada uma responde a um dos eixos
da lacuna identificada na Seção~\ref{sec:sota-lacuna}. Como o processamento roda quase inteiramente na
nuvem, a partir do catálogo público e de parâmetros versionados (Seção~\ref{sec:cost}), o procedimento
pode ser reexecutado para outros estados do Cerrado. Os sete segmentos são avaliados por um mecanismo
comum de pertinência \emph{fuzzy}/AHP e comparados em base normalizada (escore $z$), com a regra de
agregação ajustada a cada caso: média geométrica de fator limitante para os segmentos de produção e
média aritmética ponderada para a conservação (Seção~\ref{sec:aggregation}). A média geométrica se
afasta deliberadamente da agregação compensatória que predomina na área. A média aritmética retorna a
ela de forma consciente, por ser adequada a um índice de valor, e não de estado. Ambas as opções são
situadas na literatura na Seção~\ref{sec:sota-mcda}.

O agrupamento $k$-médias sobre o conjunto bruto de bandas colapsava em poucas zonas pouco
informativas, dominadas pelas bandas climáticas colineares; reduzir, ponderar por tema e aplicar a
PCA antes do agrupamento é o que torna a partição expressiva. O número de zonas, porém, é um
parâmetro escolhido, e não um achado: adota-se $k = 4$ por maximizar a reprodutibilidade da
partição em território não usado no ajuste, e não por um ótimo de geometria (Seção~\ref{sec:zoning}).
A mesma transformação climática é aplicada à linha de base e ao CMIP6. Assim, a mudança projetada é um
sinal puramente climático, o que se verifica porque fatores de mudança nulos (identidade) resultam em
$\Delta S = 0$. Os demais fatores são mantidos estáticos para isolar esse componente. A comparação com
a floresta aleatória sob amostragem espacialmente estratificada serve de verificação independente do
mapa baseado em conhecimento, e não de substituto: as divergências encontradas são tratadas como achado
sobre potencial \emph{versus} uso atual, e não como defeito oculto do modelo.

\section{Limitações}\label{sec:limitations}

Parte das limitações abaixo decorre das escolhas deliberadas de delineamento discutidas na
Seção~\ref{sec:interpretation}; outras são lacunas reais, a tratar em trabalhos futuros.

\begin{itemize}
  \item \textbf{Escopo biofísico, não econômico.} O modelo capta o potencial agroclimático, e não
    preços, posse da terra, crédito ou logística. Parte da lacuna entre potencial e uso atual é
    exatamente essa geografia econômica ausente.
  \item \textbf{Solo e terreno de produtos globais únicos.} A decomposição por fator mostra que,
    com o clima quase uniforme, solo e terreno respondem pela maior parte da discriminação espacial
    (Seção~\ref{sec:interpretation}). O atlas herda, portanto, a acurácia de dois produtos globais
    sem ajuste fino à realidade local. O solo vem do OpenLandMap a 250~m. O terreno vem do SRTM de
    1 segundo de arco ($\approx$30~m; NASA SRTM v3, com os vazios já preenchidos por fontes
    auxiliares) \citep{farr_shuttle_2007}, cuja especificação de projeto era de 16~m de erro vertical
    absoluto a 90\% de confiança; o erro estimado para a América do Sul foi de 6,2~m
    \citep[Tabela~1, p.~21]{farr_shuttle_2007}. O ruído de \emph{speckle} tende a inflar a declividade
    em escalas curtas. Essas incertezas se propagam para a elevação, a declividade e a orientação
    agregadas a 250~m por média e para o índice de umidade topográfica (interpolação bilinear).
  \item \textbf{Descompasso de escala entre clima e terreno.} As entradas climáticas
    ($\approx$4,6~km, TerraClimate) são bem mais grosseiras que as de solo e terreno (250~m), e um
    levantamento do catálogo não encontrou alternativa viável: nenhum produto reúne a normal de
    1991--2020, o conjunto completo de variáveis e uma resolução melhor
    (Seção~\ref{sec:scale-mismatch-annex}). O descompasso \emph{nominal} é, porém, enganoso. Nesta
    área, o clima varia pouco no espaço e, por isso, contribui pouco para a discriminação fina entre
    células, o que a Figura~\ref{fig:scale-mismatch} demonstra quantitativamente
    (Seção~\ref{sec:atlas-present}).
  \item \textbf{O bloco de validação não remove toda a dependência espacial.} A partição
    (Seção~\ref{sec:holdout}) usa blocos de $\approx$130~km, dimensionados pelo alcance medido do
    resíduo, mas três ressalvas limitam o que ela garante. Primeira: o bloco fica acima do alcance
    medido para soja e cana-de-açúcar, mas \emph{abaixo} do alcance das outras culturas anuais, para
    as quais não remove toda a dependência. Segunda: o alcance estimado cresce com a janela de
    observação, o que indica uma tendência regional sem alcance finito dentro do próprio território.
    As duas leituras estão quantificadas na Seção~\ref{anx:holdout}. Terceira: com cerca de quinze
    blocos por metade, as duas metades não são amostras intercambiáveis do estado. A diferença entre
    calibração e validação reflete, portanto, o território, e não otimismo do modelo (a solução é
    discutida na Seção~\ref{sec:futurework}).
  \item \textbf{O validador de produtividade é um substituto.} A NPP/GPP do MODIS mede a produção
    primária da vegetação, e não o rendimento agronômico, e a confusão entre coberturas limita seu
    uso a contexto descritivo (Seção~\ref{sec:validation-present}). AUC e Boyce seguem como
    validadores primários de presença de cultura, mas apenas onde a amostra de presença os
    sustenta. Na metade de validação, eles são fortes para a cana-de-açúcar e a soja, e
    \emph{falham} para as outras culturas anuais, cujo Boyce cai de 0,914 na metade de calibração
    para 0,313 na de validação. Para esse segmento, portanto, a tese não tem validador de presença
    utilizável.
  \item \textbf{O uso atual das outras culturas anuais não corresponde às culturas do segmento.} O
    segmento é definido como milho, algodão e sorgo (Capítulo~\ref{ch:materials}), mas a coleção do
    MapBiomas usada aqui não distingue milho nem sorgo: ambos estão diluídos na classe de
    \emph{mosaico de culturas}. A classe de uso atual com a qual o segmento é validado agrega,
    portanto, algodão, arroz e esse mosaico. É uma incompatibilidade entre o segmento modelado e o
    seu validador (não um erro do modelo) e é a explicação mais provável para o colapso do índice de
    Boyce descrito no item anterior. Todo resultado de presença das outras culturas anuais deve ser
    lido com essa ressalva; os de soja e cana-de-açúcar, que têm classe própria no MapBiomas, não são
    afetados.
  \item \textbf{A piscicultura tem uso atual quase inexistente no MapBiomas.} A aquicultura
    mapeada cobre 0,0006\% do território, cerca de 2~km\textsuperscript{2}
    (Tabela~\ref{tab:realized-composition}): uma área medida, e não uma classe ausente, mas
    pequena demais para ancorar limiares ou para validar presença. A isso se soma o fato de que
    a piscicultura é \emph{ao mesmo tempo} o único segmento
    que ganha área apta sob o clima projetado e o único cuja correlação municipal com o \emph{proxy}
    de produtividade é negativa ($\rho = -0{,}30$ na metade de validação). As duas coisas são
    consistentes, pois o validador de NPP mede vegetação terrestre e nada diz sobre aquicultura.
    Juntas, porém, significam que o segmento de maior ganho projetado é também o de menor sustentação
    empírica. A calibração da piscicultura permanece ancorada na terra genericamente
    \emph{disponível}, e não no nicho realizado, e, pela mesma razão, nenhum validador de presença
    (Boyce/AUC) é apresentado para esse segmento. Para resolver essa lacuna, seria necessário um
    mapeamento específico dos viveiros de piscicultura em GO~e~DF.
  \item \textbf{Projeção futura.} A projeção CMIP6, como qualquer projeção por fatores de
    mudança (Seção~\ref{sec:cmip6}), só pode ser sustentada indiretamente (linha de base validada e
    concordância entre GCMs). A radiação solar e a umidade do solo foram mantidas no valor histórico,
    por falta de um campo projetado e de um balanço hídrico completo, o que atenua o sinal climático
    da geração fotovoltaica e da pecuária.
\end{itemize}

\section{Trabalhos futuros}\label{sec:futurework}

A primeira frente é substituir as duas variáveis climáticas mantidas estáticas, introduzindo
radiação solar projetada e um tratamento da umidade do solo por balanço hídrico, para que a
geração fotovoltaica e a pecuária passem a ter sinal climático completo.

A segunda frente é reforçar a validação. A partição espacial por blocos (Seção~\ref{sec:holdout}) está
implementada e é a base de todas as métricas apresentadas nesta dissertação, mas deixa três pontos
abertos. O primeiro é incorporar um validador ancorado em rendimento, já que o \emph{proxy} de NPP só
sustenta uma leitura de contexto. As estatísticas da Produção Agrícola Municipal (PAM) e da Pesquisa da
Pecuária Municipal (PPM), do IBGE, fornecem rendimento por cultura e rebanho por município em série
anual. Esses dados permitiriam comparar a aptidão com a produção observada e converter a perda de área
apta projetada em uma estimativa de perda econômica. É essa a pergunta que um gestor público faz, e que
este trabalho hoje só responde em área. O segundo ponto é separar o efeito da estimação do efeito do
território. Como a maior parte das verificações avalia uma única superfície nas duas metades, a
diferença entre elas é territorial. Para isolar o ganho atribuível à estimação restrita, seria preciso
rederivar a mesma superfície a partir do domínio completo e avaliá-la novamente na metade de validação.
O terceiro ponto é substituir a partição única 50/50 por validação cruzada por blocos repetida, com
várias atribuições de blocos às metades, o que daria um intervalo em vez de um ponto. Isso importa
sobretudo para a cana-de-açúcar, cuja área cultivada é concentrada.

A terceira frente é acrescentar uma camada econômica (preços, posse da terra e infraestrutura), que
converteria parte da lacuna entre potencial e uso atual de resíduo em quantidade modelada.

A quarta frente é explorar modelos espacialmente explícitos (por exemplo, regressão com defasagem
espacial ou um termo autorregressivo condicional) para capturar efeitos de vizinhança entre
células. O mecanismo \emph{fuzzy}/AHP atual trata cada pixel de forma independente e já produz uma
interpretação espacial coerente (Seção~\ref{sec:interpretation}). Ele não modela, porém, a dependência
espacial entre células vizinhas, o que poderia refinar tanto o zoneamento quanto a leitura de incerteza
da projeção futura.

Por fim, cabe revisar os critérios de valor de conservação e a direção de sua resposta climática (perda
de aptidão \emph{versus} prioridade crescente de proteção). Isso pode ser feito no espírito do
planejamento multicritério de compromissos entre restauração, serviços ecossistêmicos e aptidão
agrícola já aplicado ao Cerrado \citep{schuler_spatial_2022}, estendendo-o à dimensão climática
(Seção~\ref{sec:sota-casos}).

\section{Conclusão}\label{sec:conclusion}

Esta dissertação mapeou o potencial agroclimático de sete segmentos de uso da terra em Goiás e
no Distrito Federal, particionou o território em zonas agroambientais e projetou como esse
potencial se desloca com a mudança do clima até meados do século, com um método aberto e
reprodutível na nuvem.

As três questões de pesquisa foram respondidas. Um mecanismo \emph{fuzzy}/AHP produz superfícies de
aptidão do presente para os sete segmentos. Um agrupamento não supervisionado divide o território em
quatro zonas que, construídas sem nenhuma variável de cobertura, separam de forma independente a
produtividade observada, ainda que, em duas delas, que somam três quartos da área, a vocação
comparativa seja indistinta entre os dois segmentos mais bem colocados. A projeção CMIP6 mostra, até
2070, uma mudança modesta, mas com sinal consistente entre os modelos. A geração fotovoltaica é o
segmento mais exposto, seguida da conservação e das outras culturas anuais; soja, cana-de-açúcar e
pecuária permanecem estáveis, e a piscicultura é a única a ganhar área apta (QP3). Uma lacuna entre
potencial e uso atual, validada espacialmente e concentrada em pastagem existente, aponta para a
intensificação como direção mais provável de mudança de uso da terra na região.

De forma mais geral, o trabalho mostra que uma avaliação de terras multissegmento, sensível ao
clima e comparada com o uso atual, pode ser montada inteiramente a partir de dados públicos. O
atlas resultante indica onde cada segmento tem maior aptidão em GO~e~DF, onde esse potencial ainda não
é aproveitado e como a mudança do clima deve redesenhá-lo até meados do século.
